\addchap{Preface}

This book began with a list of textbooks that should exist but do not. One entry observed that neural networks are singular statistical models: many different weights give the same function, the Fisher information is degenerate, and the asymptotic theory behind the Laplace approximation, the Bayesian information criterion, and Akaike's criterion does not apply to them. The entry noted that the correct theory exists. Sumio Watanabe built it from 1999 on, and it replaces the familiar half the number of parameters by a single invariant of the loss, the learning coefficient. It has recently been used to measure how the complexity of a network changes during training. But Watanabe's monograph of 2009 asks its readers to know resolution of singularities before they begin, and the theory has a reputation for being hard to learn. The entry proposed a textbook for the rest of us: one family of toy models, starting from the model $y=abx$ with two parameters, on which every concept is computed by hand; Hironaka's theorem on resolution of singularities used as a black box, stated precisely and never proved; and a last part that meets the empirical results of deep learning.

The resulting book has fifteen chapters in four parts. Part I sets the classical theory against the running toy models and computes the free energy of the simplest of them completely, which turns out to contain most of the theory in miniature. Part II is about the learning coefficient as a geometric quantity: the rate at which the volume of nearly true parameters shrinks, computed by blowing up the parameter space or by reading a Newton polygon, with a catalogue of the values that are known. Part III is about Watanabe's theorems on the free energy, the generalization error, the criteria WAIC and WBIC, and the posterior's response to temperature: it proves the central result of each completely for a class of models where the computation can be carried through, states the general theorem with its source, and shows how a posterior moves between singularities as the data grow. Part IV asks what all this says about trained networks: the local learning coefficient, its changes during training, and what singular learning theory does not explain. Each chapter states what the regular theory says about its subject, computes the singular answer by hand on a toy model, proves its central theorem completely, in general or for such a class of models, measures the result in a short script, records it on a ledger card, and closes by saying where the theory stops.

The intended readers are researchers and graduate students in machine learning and statistics who have met the claim that neural networks are singular and want the theory behind it. The prerequisites are multivariable calculus, linear algebra, and probability at the level of a first graduate course; no algebraic geometry is assumed, and none beyond the statement of one theorem is used. Each chapter defines what it needs, and an appendix collects the shared tools, with proofs of those whose proofs are short.

Two cautions belong at the front. First, singular learning theory is a theory of exact Bayesian posteriors as the number of examples grows with the model fixed. Deep learning trains one network by stochastic gradient descent, often with more parameters than examples. The book tries to say exactly which results transfer, which transfer under assumptions that can be checked, and which do not transfer at all; the last chapter is devoted to the gap. Second, much of the work that connects the theory to deep learning is young. Several of the results discussed are preprints from 2023 to 2026, and most of the experiments behind them are small. The chapters date such claims and say when a result has not been refereed.

\paragraph{A note on authorship.} This book was written by Claude, an AI model made by Anthropic, and conceived and edited by Xiaoyu Li, who proposed the subject, its audience, the gap it fills, and its approach through one family of toy models. Every theorem is stated with its hypotheses and is proved, sketched, or attributed to the source that proves it; several results and proofs are the book's own and are marked as such. Every claim attributed to a paper was checked against the paper or a reliable account of it wherever these could be reached, every number from the book's scripts comes from code that the reader can run, and the colophon describes how the book was made. Readers who find an error are invited to report it, with the computation that fails.
